\subsection{Visualização Tática e Segmentação Operacional}

A materialização matemática da Consciência Situacional é concretizada através da camada de apresentação no Apache Superset. Para mitigar a carga cognitiva de operadores em salas de situação, a construção dos \textit{dashboards} seguiu o princípio da segmentação tática, dividindo o monitoramento em domínios operacionais distintos:

\begin{itemize}
    \item \textbf{Visão Macro (Biomas e Sociedade):} Focada na análise exploratória geral, disponibilizando volumetria histórica, tendências temporais e mapas de calor consolidados em níveis estaduais e por biomas.
    \item \textbf{Vulnerabilidade Humanitária e Sanitária:} Interface direcionada à mitigação de crises diretas, evidenciando as métricas de pressão exercida pelo fogo sobre territórios sensíveis (Terras Indígenas) e sobre a rede de suporte crítico de saúde (Hospitais e Unidades de Pronto Atendimento).
    \item \textbf{Risco Estrutural em Redes Complexas:} Painel projetado para a equipe de infraestrutura, traduzindo visualmente os cálculos da Teoria dos Grafos. A interface destaca os nós da malha elétrica com base em sua Centralidade de Intermediação e estresse térmico, permitindo a antecipação de falhas em cascata.
\end{itemize}

Essa separação garante que cada equipe técnica consuma a parcela dos dados pertinente à sua esfera de atuação, amparada pelas mesmas Tabelas Fato padronizadas em alta performance.
